\chapter{REFORMANDO LAS REGLAS DE JUEGO}\label{chap:normas}

Ante un panorama caracterizado por un elevado nivel de endeudamiento, déficits persistentes y una presión creciente de diverso tipo de transferencias, se hace imperativo repensar el modelo fiscal del país para garantizar su sostenibilidad a largo plazo. Este capítulo formula algunas propuestas para enfrentar estos desafíos. Sin embargo, se necesita ir más allá de la reiterada necesidad de una reforma tributaria estructural y de la búsqueda de una mayor eficiencia del gasto. Se requiere también un replanteamiento del marco constitucional, legal y administrativo bajo el cual se formula y ejecuta la política fiscal.

La activación de la cláusula de escape en junio de 2025 constituye la expresión más clara de las limitaciones del actual marco fiscal. Ante la imposibilidad de conciliar las restricciones de la regla fiscal con las obligaciones de gasto legalmente establecidas, el Gobierno optó por suspender temporalmente los límites al déficit fiscal. Esta decisión deja en evidencia que el marco institucional colombiano carece de mecanismos eficaces para enfrentar situaciones en las que los ingresos proyectados no alcanzan para financiar el gasto legalmente exigido.

Además del ejercicio de consolidación fiscal que requiere ajustes en los ingresos y el gasto para Colombia en los próximos años, es urgente un cambio en el marco normativo y administrativo que rige la formulación del presupuesto. Este cambio es necesario para superar una práctica insostenible: la formulación de presupuestos que no reflejan de manera realista la capacidad de financiación del Estado. Como lo evidenció la activación de la cláusula de escape en 2025, mantener esta dinámica solo posterga los ajustes necesarios y debilita la credibilidad del marco fiscal.

El problema de los presupuestos inflados no obedece a un exceso de optimismo ni a decisiones políticas coyunturales. Responde a una tensión estructural: el Ministerio de Hacienda está legalmente obligado a incluir en el \gls{pgn} todas las erogaciones derivadas de mandatos constitucionales y legales, sin contar necesariamente con ingresos suficientes para financiarlas. En ausencia de un mecanismo institucional que reconozca y resuelva este desajuste, el Gobierno ha recurrido a la sobreestimación sistemática de ingresos o, como en 2025, a la activación de la cláusula de escape para evitar presentar un presupuesto abiertamente incompatible con la regla fiscal. Las propuestas que se presentan a continuación buscan corregir este desequilibrio, estableciendo que, si luego de surtido sin éxito el trámite de una Ley de Financiamiento, las proyecciones de ingresos —validadas de forma técnica e independiente— no alcanzan para cubrir las obligaciones legales de gasto —también validadas de manera independiente—, el Ejecutivo declare la emergencia económica y, por lo tanto, quede facultado para decretar los impuestos adicionales requeridos para cumplir con la regla fiscal. Con ello se busca una presupuestación realista, coherente con la regla fiscal, pero que permita dar cabida al reconocimiento pleno de las obligaciones de gasto de origen legal.

\section{Oficina de Asistencia Técnica Presupuestal}

Una pieza clave para el fortalecimiento del marco fiscal institucional en Colombia es la implementación de la Oficina de Asistencia Técnica Presupuestal (\gls{oatp}), adscrita al Congreso de la República y creada por la \href{https://www.funcionpublica.gov.co/eva/gestornormativo/norma.php?i=98370}{Ley 1985 de 2019}. Esta entidad está pensada para dotar al Legislativo de capacidades técnicas independientes para analizar, proyectar y debatir de forma rigurosa las políticas fiscales y presupuestales propuestas por el Ejecutivo.

Actualmente, el Congreso aprueba el Presupuesto General de la Nación y tiene competencias en materia fiscal, pero enfrenta limitaciones evidentes en su capacidad técnica para evaluar con profundidad la información que recibe. La implementación de la \gls{oatp} permitiría reducir esa asimetría entre el Ejecutivo — el cual cuenta con el respaldo técnico del Ministerio de Hacienda— y el Legislativo, que carece de un cuerpo permanente y especializado que le proporcione insumos analíticos sólidos.

La experiencia internacional muestra que este tipo de oficinas, como el Congressional Budget Office (CBO) en Estados Unidos o el Parliamentary Budget Officer (PBO) en Canadá, pueden jugar un papel determinante en el fortalecimiento de la transparencia fiscal y del control democrático de los recursos públicos. Estas entidades permiten que las decisiones presupuestales se tomen con base en evidencia, y no únicamente a partir de las cifras y narrativas presentadas por el Ejecutivo.

Además, a diferencia de otras instituciones fiscales independientes, los conceptos emitidos por una \gls{oatp} pueden adquirir un carácter vinculante en el trámite legislativo, en tanto se integran al proceso de deliberación y aprobación de leyes. Esto contrasta con el rol de instituciones como el Comité Autónomo de la Regla Fiscal (\gls{carf}), cuyos conceptos, aunque técnicamente valiosos, difícilmente pueden llegar a ser vinculantes, pues su mandato se limita a la evaluación del cumplimiento de la regla fiscal.

Si bien el \gls{carf} constituye un avance importante en la arquitectura institucional del marco fiscal colombiano, no puede reemplazar la función que cumpliría la \gls{oatp}. En primer lugar, su naturaleza es distinta: el \gls{carf} está diseñado como un órgano de vigilancia y evaluación de la sostenibilidad fiscal y del cumplimiento de la regla fiscal, no como un ente asesor del Congreso. Su relación institucional es más cercana al Ejecutivo, en tanto evalúa sus proyecciones y políticas, y no tiene un mandato para interactuar activamente con el Legislativo.

En segundo lugar, el \gls{carf} cuenta con capacidades limitadas y un enfoque eminentemente macrofiscal. Sus análisis se concentran en agregados como el balance fiscal, la trayectoria de deuda y la consistencia de las proyecciones oficiales, sin que ello implique una capacidad operativa para desagregar el análisis a nivel del gasto sectorial o los efectos distributivos del gasto y los impuestos, como sí lo podría hacer la \gls{oatp}.

La implementación de la \gls{oatp} requiere ajustes a la ley que la creó. El concurso de méritos para seleccionar al director y a los subdirectores es un obstáculo operativo, y el período de apenas dos años para el director limita la continuidad técnica. Una reforma legal que simplifique el proceso de selección y amplíe el período a cuatro años daría mayor estabilidad a la entidad y facilitaría su puesta en marcha.

En síntesis, el diseño de una institucionalidad fiscal robusta en Colombia requiere entender que la \gls{oatp} y el \gls{carf} cumplen funciones complementarias pero no sustituibles. Mientras el \gls{carf} evalúa el desempeño macrofiscal del Gobierno, la \gls{oatp} acompañaría el proceso legislativo con herramientas de análisis independientes, ofreciendo al Congreso una capacidad real de escrutinio y deliberación informada sobre las decisiones fiscales del país.

\section{Validación de proyecciones de ingresos y gastos}

La validación de la meta de recaudo por parte del \gls{carf} debe ser vinculante. Esto con el fin de garantizar que el \gls{pgn} sea formulado en estricta coherencia con la regla fiscal. Presentar un \gls{pgn} sin la validación previa de la meta de recaudo por parte del \gls{carf} implicaría la violación al mandato legal de coherencia del \gls{pgn} con la regla fiscal.  

En línea con lo anterior, resulta conveniente establecer una distinción explícita entre la meta de recaudo institucional de la \gls{dian} y la proyección de ingresos tributarios que se incorpora en el \gls{pgn}. La meta de la \gls{dian} puede —y debe— ser ambiciosa, con el objetivo de incentivar la gestión administrativa de la \gls{dian}. No obstante, utilizar dicha meta como base para la programación presupuestal puede introduce un riesgo para la sostenibilidad fiscal, ya que amplía artificialmente el espacio de gasto y genera el riesgo de incumplir la regla fiscal en la medida en que la meta de la \gls{dian} no se materialice y no se alcancen a implementar los respectivos ajustes de gasto. Para evitar este riesgo, la proyección de ingresos en el \gls{pgn} debe corresponder a un escenario realista y validado por el \gls{carf}, mientras que la meta de la \gls{dian} debe funcionar como un instrumento de gestión y evaluación de desempeño, sin trasladar sus supuestos optimistas a la programación fiscal de la Nación.

De forma complementaria, la estimación del monto de las obligaciones de gasto de origen legal y constitucional —incluyendo aquellas derivadas del \gls{sgp}, las pensiones, la salud y demás mandatos legales— también debe ser objeto de validación independiente. Esta validación, a cargo de la \gls{oatp} permitiría garantizar que las proyecciones de gasto obligatorio no sean sobrestimadas. De este modo, el monto de las inflexibilidades legales se convertiría en un insumo obligatorio en la formulación del \gls{pgn}, y su validación por parte de una institución fiscal independiente sería condición necesaria para certificar la viabilidad del presupuesto tanto en el marco de la regla fiscal como en el marco de las inflexibilidades legales del gasto.

\section{Consistencia entre regla fiscal y las obligaciones legales de gasto} 

No obstante, para que el \gls{pgn} pueda formularse de manera coherente con la regla fiscal al mismo tiempo garantice el cumplimiento de los compromisos legales y constitucionales de gasto, es fundamental que cualquier insuficiencia de recursos sea identificada y comunicada oportunamente por el Ministerio de Hacienda. En estos casos, la presentación del presupuesto debe ser acompañado de una Ley de Financiamiento siempre que los ingresos no sean suficientes para cubrir las inflexibilidades legales. En este caso, la proyección de ingresos debe estar validada por el \gls{carf} mientras la estimación de las inflexibilidades legales debe estar validada por la \gls{oatp}.

%En caso de que dicha Ley de Financiamiento no sea aprobada o los recursos resulten insuficientes para garantizar el cumplimiento de las inflexibilidades legales, deberá declararse el estado de cosas inconstitucional, dado que en este escenario el Estado no contaría con los recursos para hacer cumplir los derechos constitucionales. Esto facultaría al Gobierno para decretar impuestos que permitan cubrir dichas obligaciones, asegurando así la sostenibilidad fiscal y el cumplimiento de los compromisos establecidos en la Constitución y la legislación vigente. La Figura \ref{fig:flujoPGN} resume esta propuesta para un ciclo de formulación del presupuesto consistente con la regla fiscal sin correr el riesgo de incumplir con las obligaciones de gasto de origen legal. 

En caso de que dicha Ley de Financiamiento no sea aprobada o los recursos aprobados resulten siendo insuficientes para garantizar simultáneamente el cumplimiento de las obligaciones legales de gasto y de la regla fiscal, el Gobierno debería estar obligado a declarar la emergencia económica y por ende quedar facultado para decretar los impuestos que permitan cubrir dichas obligaciones. De esta manera se preservaría tanto la sostenibilidad fiscal como el cumplimiento de los compromisos establecidos en la Constitución y la legislación vigente. La Figura \ref{fig:flujoPGN} resume esta propuesta para un ciclo de formulación del presupuesto consistente con la regla fiscal sin correr el riesgo de incumplir con las obligaciones de gasto de origen legal.

\begin{figure}[h!]
    \centering
    \caption{Propuesta para el ciclo de aprobación del PGN}
    \label{fig:flujoPGN}
    \vspace{0.3cm}
\begin{tikzpicture}[
    node distance = 0.8cm, % menor distancia vertical
    block/.style = {
        rectangle,
        draw,
        text width=12cm, % más ancho
        minimum height=0.6cm, % más bajo
        align=center
    },
    arrow/.style = {-Stealth, thick}
]

% Nodos
\node[block] (1) {1. CARF valida proyección de ingresos y OATP valida inflexibilidades legales de gasto};
\node[block, below=of 1] (2) {2. Hacienda presenta PGN con todas las obligaciones legales de gasto};
\node[block, below=of 2] (3) {3. Si los ingresos proyectados son insuficientes, Hacienda anuncia riesgo de incumplir con la regla fiscal.};
\node[block, below=of 3] (4) {4. Se presenta Ley de Financiamiento en caso de que proyecciones de ingresos y obligaciones legales de gasto sean inconsistentes con la regla fiscal};
%\node[block, below=of 4] (5) {5. Si la Ley de Financiamiento no es aprobada, el Gobierno decreta los impuestos apenas suficientes para cumplir con las inflexibilidades legales.};
\node[block, below=of 4] (5) {5. Si la Ley de Financiamiento no es aprobada se declara la emergencia económica};
\node[block, below=of 5] (6) {6. El Gobierno decreta los impuestos apenas suficientes para cumplir con las inflexibilidades legales y con la regla fiscal.};

% Flechas
\draw[arrow] (1) -- (2);
\draw[arrow] (2) -- (3);
\draw[arrow] (3) -- (4);
\draw[arrow] (4) -- (5);
\draw[arrow] (5) -- (6);
\end{tikzpicture}
\end{figure}


Aunque  en esta propuesta se restringen los márgenes de maniobra del Gobierno en la etapa de formulación presupuestal —al exigir que tanto las proyecciones de ingresos como las obligaciones legales de gasto sean validadas de forma independiente y vinculante—, se le reconoce una facultad excepcional en una etapa siguiente, al permitirle decretar impuestos cuando el Congreso no apruebe una Ley de Financiamiento que garantice la coherencia fiscal. En otras palabras, se pierde discrecionalidad en la fase inicial, pero se gana capacidad decisoria ante escenarios de bloqueo político, lo que refuerza la credibilidad de la regla fiscal como ancla institucional y asegura el cumplimiento de los compromisos del Estado.

\section{Efectos de la no aprobación del presupuesto}

En la actualidad, el artículo 348 de la Constitución establece dos situaciones distintas para asegurar la continuidad del funcionamiento del Estado ante la ausencia de una ley anual de presupuesto. Por un lado, si el Congreso no aprueba el proyecto de presupuesto dentro del término constitucional, entra a regir automáticamente el proyecto presentado por el Gobierno, tal como fue radicado. Por otro lado, si el Gobierno no presenta el presupuesto dentro del plazo establecido, entra a regir el presupuesto del año anterior, y en ese caso el Ejecutivo queda facultado para reducir gastos y, en consecuencia, suprimir o refundir empleos.

Cuando el Congreso no aprueba el presupuesto —como ocurrió en 2024 con el presupuesto para la vigencia de 2025— el Gobierno se ve beneficiado con la entrada en vigencia automática de su proyecto original, sin modificaciones, a pesar de las reservas expresadas por el Legislativo frente a las proyecciones de ingresos incluidas en dicho proyecto. Esta experiencia pone de manifiesto cómo la versión actual de la norma debilita el papel del Congreso como órgano de control político y presupuestal.

Se propone, por tanto, modificar el artículo 348 para igualar las consecuencias de ambas omisiones. Es decir, que si el Congreso no expide el presupuesto en el plazo constitucional, no entre a regir automáticamente el proyecto del Gobierno, sino que se reconduce el presupuesto del año anterior, aplicándose en ese caso las mismas reglas previstas cuando el presupuesto no ha sido presentado oportunamente. Esto implica que el Gobierno podría reducir gastos y reestructurar empleos si las condiciones fiscales así lo exigen.

%\section{Contabilización de los compromisos como gasto fiscal}

%Como se ha discutido, el crecimiento de los compromisos adquiridos en ejercicios anteriores que no pudieron ser obligados a tiempo plantea un desafío significativo para la gestión fiscal. Este fenómeno ha generado una acumulación de reservas (compromisos no obligados) que compromete tanto la liquidez como la sostenibilidad del presupuesto público. Esta problemática ha surgido debido a la estrategia de inflar la meta de recaudo para abrir espacio presupuestal para más apropiaciones. En la medida en que esas metas de recaudo no se han materializado, las restricciones en la disponibilidad de caja han hecho que los compromisos no alcancen a obligarse para el año en curso, aplazándose y acumulándose para el año siguiente.

%Para resolver esta problemática, el reconocimiento de los gastos debería hacerse en el momento en que se comprometen, en lugar de cuando se obligan, lo que permitiría una alineación más inmediata del presupuesto con la regla fiscal. Naturalmente, las pérdidas de apropiación —es decir, recursos presupuestados que no logran comprometerse dentro del ejercicio fiscal— se seguirían convirtiendo en ahorros fiscales que excederían las metas de la regla fiscal. En la medida en que los compromisos no alcancen a obligarse, estos ahorros contribuirían a aumentar el superávit primario, proporcionando mayor capacidad para reducir la deuda pública y mejorar la sostenibilidad fiscal a largo plazo. 

%\section{Reconducción automática del presupuesto anterior como mecanismo de contingencia}

%El diseño actual del artículo 348 de la Constitución, que establece que el presupuesto presentado por el Gobierno entra a regir automáticamente si no es aprobado por el Congreso antes del 20 de octubre, otorga al Ejecutivo un poder excesivo en la negociación presupuestal y debilita el principio de aprobación parlamentaria. Una alternativa institucional más equilibrada sería establecer que, en ausencia de aprobación oportuna, quede vigente el presupuesto del año anterior, con los ajustes necesarios para garantizar el funcionamiento básico del Estado. Este mecanismo de reconducción presupuestaria se aplica en países como España, Francia e Italia, y contribuye a evitar tanto el cierre del gobierno como la imposición automática del proyecto del Ejecutivo. De este modo, se protege la continuidad del gasto público esencial, se preserva el control legislativo sobre el presupuesto y se evita que el Gobierno tenga incentivos para forzar un desacuerdo con el Congreso con el fin de imponer su propuesta sin modificaciones.

%\section{Fortalecimiento del aval fiscal}

%Aunque en Colombia existe el aval fiscal, su implementación ha perdido efectividad, reduciéndose a un visto bueno por parte del Ministerio de Hacienda, el cual es políticamente difícil de negar cuando la iniciativa legislativa proviene del Ejecutivo. Se debe exigir que toda iniciativa legislativa que incremente el gasto o reduzca ingresos venga acompañada de una fuente de financiación explícita. Es decir, que el proyecto identifique de forma concreta qué partidas presupuestales se reducirán, qué impuestos adicionales se crearán o qué tarifas existentes se incrementarán para cubrir su impacto fiscal. Esta exigencia fortalecería la responsabilidad fiscal al obligar a los proponentes a enfrentar los costos políticos de sus iniciativas.

%Aunque en Colombia existe el aval fiscal, su implementación ha perdido efectividad. En la práctica, basta con que el Ministerio de Hacienda declare que una iniciativa es “compatible con el Marco Fiscal de Mediano Plazo” para que avance en el Congreso, incluso si existen dudas sobre su sostenibilidad. Además, cuando la iniciativa proviene del propio Ejecutivo, resulta políticamente difícil que Hacienda niegue el aval, lo que reduce su valor como herramienta de control fiscal.

%Para corregir este desequilibrio, se propone establecer un esquema de control cruzado que distinga el origen de la iniciativa. Cuando el proyecto provenga del Congreso y tenga impacto fiscal, el aval debe ser emitido por la \gls{oatp}, con base en una estimación del costo y la fuente de financiación. El Ministerio de Hacienda, en este caso, validaría si la propuesta es compatible con el Marco Fiscal de Mediano Plazo. Si, por el contrario, la iniciativa proviene del Ejecutivo, el Ministerio de Hacienda emitiría el aval, pero este requeriría la validación de la \gls{oatp} antes de continuar su trámite.

%Adicionalmente, se debe exigir que toda iniciativa legislativa que incremente el gasto o reduzca ingresos venga acompañada de una fuente de financiación explícita. Es decir, que el proyecto identifique de forma concreta qué partidas presupuestales se reducirán, qué impuestos adicionales se crearán o qué tarifas existentes se incrementarán para cubrir su impacto fiscal. Esta exigencia fortalecería la responsabilidad fiscal al obligar a los proponentes a enfrentar los costos políticos de sus iniciativas.

\section{Unificación de la formulación del \gls{pgn}}

Tal como lo recomendó la Comisión de Gasto e Inversión Pública (\cite{comisiongasto2017}), es necesario unificar el proceso presupuestal de funcionamiento e inversión en una sola entidad, preferiblemente el Ministerio de Hacienda. Esto permitiría resolver el conflicto institucional actual, en el que el Departamento Nacional de Planeación (\gls{dnp}) gestiona el presupuesto de inversión, mientras que Hacienda tiene el mandato de garantizar la sostenibilidad fiscal.

Las inflexibilidades presupuestales hacen que, al momento de ajustar el gasto, la inversión sea el rubro más susceptible a recortes. Sin embargo, el \gls{dnp} suele mostrar una resistencia significativa a estas decisiones, lo que complica la labor del Ministerio de Hacienda para cumplir con las metas fiscales establecidas.

%Centralizar el manejo del presupuesto en Hacienda permitiría una mayor coherencia en las políticas de gasto público y fortalecería la capacidad de reacción frente a ajustes fiscales. Esto eliminaría duplicidades y facilitaría una toma de decisiones más alineada con las prioridades macroeconómicas del país.

Centralizar el manejo del presupuesto en Hacienda permitiría una mayor coherencia en las políticas de gasto público y fortalecería la capacidad de reacción frente a ajustes fiscales. No obstante, esta reforma debe preservar y aprovechar la capacidad técnica que el DNP ha desarrollado en la formulación, seguimiento y evaluación de proyectos de inversión pública. Esto podría implicar el traslado de personal y capacidades técnicas del DNP al Ministerio de Hacienda, con el fin de integrar dichas competencias en un proceso presupuestal unificado y más consistente con los objetivos fiscales y macroeconómicos del país, sin desmantelar el capital institucional construido en torno a la inversión pública.

%\section{Trimestralización de las apropiaciones}

%La situación actual evidencia las consecuencias de la ausencia de herramientas efectivas para ajustar las políticas fiscales en tiempo real. Como solución, se debe implementar que las apropiaciones presupuestales estén trimestralizadas y que el Ministerio de Hacienda tenga la potestad de hacer ajustes (mediante actos administrativos y no mediante decretos gubernamentales) sobre las apropiaciones trimestrales basadas en la ejecución del trimestre anterior. Aunque estos ajustes no estarán exentos de controversias al interior del gobierno, resultarían más efectivos y transparentes que el manejo de caja como herramienta de asignación.

%Los problemas fiscales que se sufren en la actualidad también son consecuencia de la falta de capacidad por parte del Ministerio de Hacienda de ajustar y reasignar las apropiaciones en el transcurso de cada año. Por lo tanto, se debe permitir al Ministerio realizar ajustes trimestrales a las apropiaciones presupuestarias. Esto  dotaría al presupuesto de mayor flexibilidad y adaptabilidad, sin comprometer los principios de transparencia y rendición de cuentas. Esta medida permitiría responder de manera más efectiva a fluctuaciones en los ingresos fiscales o cambios en las prioridades de la Nación, asegurando que los recursos públicos se destinen a donde son más necesarios durante el año fiscal y colaborando con las metas fiscales en escenarios de caída de los ingresos fiscales.

%Para implementar este mecanismo, sería necesario establecer un marco normativo que regule los ajustes trimestrales bajo criterios claros, como subejecución presupuestaria, ingresos excedentes o necesidades emergentes. Cada ajuste debería estar acompañado de informes públicos y la aprobación por parte de organismos de control, garantizando la transparencia y evitando abusos. %Este enfoque no solo reduciría la subejecución presupuestaria y optimizaría el uso de recursos, sino que también fortalecería la capacidad del Gobierno para garantizar la estabilidad fiscal.

\section{Medición y seguimiento a la ejecución presupuestal}

La ejecución presupuestal en Colombia se mide tradicionalmente como el nivel de obligaciones contraídas en relación con las apropiaciones aprobadas en el presupuesto. Sin embargo, este indicador se ha convertido en la principal métrica de desempeño, desplazando la atención del seguimiento a los resultados y objetivos específicos de cada programa de gasto.

Es fundamental que el país avance hacia un modelo de presupuestación basada en resultados, como lo plantea el Documento CONPES 4083 (\cite{conpes4083}). Este enfoque busca vincular explícitamente los recursos asignados con los resultados e impactos esperados, utilizando indicadores de desempeño que permitan evaluar la eficacia y eficiencia del gasto público y mejorar la rendición de cuentas.

Mientras se consolida la implementación de este modelo, resulta necesario introducir mejoras inmediatas en la medición de la ejecución presupuestal. En particular, se debe excluir de los indicadores de ejecución las transferencias de recursos a fiducias y patrimonios autónomos. Aunque estas transferencias constituyen un compromiso formal, no garantizan que los recursos se ejecuten de manera efectiva y oportuna en los fines previstos. De hecho, muchas entidades las utilizan estratégicamente para inflar sus cifras de ejecución y evitar recortes o reasignaciones presupuestales. Corregir esta práctica permitiría al Ministerio de Hacienda realizar ajustes y reasignaciones con mayor precisión y oportunidad, al tiempo que aumentaría significativamente la transparencia y la calidad de la información sobre la ejecución del gasto.

%Dicho esto, la ejecución presupuesta debe ser reevaluada. La ejecución presupuestal se mide como el nivel de obligaciones contraídas en relación con las apropiaciones aprobadas en el presupuesto. Sin embargo, el país ha adoptado este indicador como la principal métrica de desempeño, en lugar de enfocarse en establecer y hacer seguimiento a metas vinculadas con los objetivos específicos de cada programa de gasto.

%Mientras se avanza en la implementación de una presupuestación por programas y una evaluación basada en resultados, es necesario descontar de los indicadores de ejecución presupuestal a las transferencias de recursos a fiducias y patrimonios autónomos. Estas transferencias, aunque representan un compromiso formal, no garantizan que los recursos se utilicen de manera efectiva y oportuna en los propósitos para los que fueron asignados. Lamentablemente, las entidades del Gobierno suelen emplear esta práctica para inflar sus indicadores de ejecución, protegiéndose así de posibles ajustes o reasignaciones de gasto. Implementar esta medida fortalecería la capacidad del Ministerio de Hacienda para realizar ajustes y reasignaciones de manera más efectiva, al tiempo que aumentaría significativamente la transparencia y la precisión en la medición de la ejecución presupuestal.

\section{Reformas a la figura de pérdidas de apropiación}

El principio de anualidad presupuestal establece que las apropiaciones de gasto solo tienen validez durante el año fiscal en que fueron aprobadas. Bajo esta lógica, las apropiaciones presupuestadas que no se comprometen durante la vigencia pierden su validez, lo que se conoce como pérdidas de apropiación. Sin embargo, esta figura genera incentivos perversos que presionan a las entidades a ejecutar recursos apresuradamente hacia el final del año con el fin de evitar su pérdida. Este comportamiento no solo compromete la eficiencia y calidad del gasto público, sino que también agrava la presión sobre el cierre fiscal y sobre la disponibilidad de caja.

En contraste, países como Estados Unidos clasifican las apropiaciones según su duración —anuales, multianuales o indefinidas— y permiten la ejecución de pagos más allá del año fiscal. Australia, por su parte, permite que las apropiaciones se ejecuten en ejercicios futuros sin necesidad de reautorización. El Reino Unido, por su parte, ha implementado un esquema mediante el cual los ministerios pueden trasladar una parte predeterminada de los recursos no ejecutados al siguiente ejercicio fiscal, previa aprobación por parte del Tesoro. 

Colombia no debería excluir la posibilidad de eliminar la figura de pérdidas de apropiación, o al menos reformarla, permitiendo el traslado parcial de saldos no comprometidos bajo reglas claras de programación plurianual. 

%El principio de anualidad presupuestal establece que las apropiaciones de gasto solo tienen validez durante el año fiscal en que fueron aprobadas. Bajo esta lógica, las apropiaciones presupuestadas que no se comprometen durante la vigencia fiscal pierden su validez, lo que se conoce como pérdidas de apropiación. Sin embargo, esta figura genera incentivos perversos que presionan a las entidades a gastar apresuradamente hacia el final del año con el fin de evitar la pérdida de recursos. Este comportamiento no solo compromete la eficiencia y calidad del gasto público, sino que también agrava la presión sobre el cierre fiscal y sobre la disponibilidad de caja. En contraste, países como Estados Unidos clasifican las apropiaciones según su duración —anuales, multianuales o indefinidas— y permiten la ejecución de pagos más allá del año fiscal, mientras que en Australia las apropiaciones no expiran automáticamente y pueden ejecutarse en ejercicios futuros sin necesidad de reautorización. Colombia no debería excluir la posibilidad de eliminar la figura de pérdidas de apropiación, o al menos reformarla, permitiendo el traslado parcial de saldos no comprometidos bajo reglas claras de programación plurianual, como ocurre en muchos países con marcos fiscales más modernos.

%\section{Suspensión de las Transacciones de Única Vez}

%Algunos cambios a la regla fiscal deben ser considerados, El primero es suspender las Transacciones de Única Vez (\gls{tuv}), las cuales consisten en ingresos o gastos excepcionales que no se repiten regularmente y, por lo tanto, no forman parte de los ingresos o egresos estructurales del Gobierno. Estas transacciones pueden incluir, por ejemplo, la venta de activos estatales o gastos excepcionales derivados de desastres naturales. No obstante, aunque la \href{https://www.funcionpublica.gov.co/eva/gestornormativo/norma.php?i=170902}{Ley 2155 de 2021} especifica los criterios para que una transacción sea considerada como una \gls{tuv}, es el Consejo de Política Fiscal (CONFIS) quien define si una transacción es una \gls{tuv}, previo concepto no vinculante del Comité Autónomo de la Regla Fiscal (\gls{carf}). Esto introduce un sesgo en el sentido de que es el Gobierno a través del CONFIS quien selecciona qué transacciones deben ser consideradas como una \gls{tuv}. Naturalmente, el Gobierno selecciona las \gls{tuv} que abren espacio fiscal (caída de ingresos o aumentos de gasto extraordinarios) e ignora las que lo cierran (aumento de ingresos o caídas de gasto extraordinarios). Adicionalmente, si el PGN se formula de manera \textit{ex-ante} de una manera consistente con la regla fiscal y las proyecciones de ingresos y gastos están validadas de manera técnica e independiente, las variaciones imprevistas de ingresos no deberían generar ni ajustes sobre el gasto ni controversias sobre el cumplimiento de la regla fiscal.

\section{Revisión del tratamiento de las Transacciones de Única Vez}

Algunos cambios a la regla fiscal deben ser considerados. Uno de ellos es reformar el tratamiento actual de las Transacciones de Única Vez (\gls{tuv}), entendidas como ingresos o gastos excepcionales que no se repiten regularmente y, por lo tanto, no forman parte de los ingresos o egresos estructurales del Gobierno. Estas transacciones pueden incluir, por ejemplo, la venta de activos estatales o gastos extraordinarios derivados de desastres naturales.

Aunque la \href{https://www.funcionpublica.gov.co/eva/gestornormativo/norma.php?i=170902}{Ley 2155 de 2021} especifica los criterios para que una transacción sea considerada como una \gls{tuv}, en la práctica es el Consejo de Política Fiscal (CONFIS) quien toma la decisión final, previo concepto no vinculante del Comité Autónomo de la Regla Fiscal (\gls{carf}). Esto introduce un sesgo sistemático: el Gobierno tiende a reconocer como \gls{tuv} aquellas transacciones que amplían el espacio fiscal (caídas de ingresos o aumentos de gasto extraordinarios), pero omite aplicar el mismo tratamiento a transacciones que lo reducirían (como aumentos inesperados de ingreso o reducciones de gasto no recurrentes).

Para corregir esta asimetría y fortalecer la credibilidad del marco fiscal, se propone que la clasificación de una transacción como \gls{tuv} sea definida exclusivamente por el \gls{carf}, con base en criterios técnicos y bajo un procedimiento público y transparente. El Gobierno podría seguir presentando propuestas, pero la decisión final recaería en un órgano independiente. Esta reforma preserva la función técnica de las \gls{tuv} como mecanismo para aislar los efectos transitorios en el balance fiscal, pero elimina el margen de discrecionalidad política que ha erosionado su legitimidad.

\section{Activación de la cláusula de escape}

La activación de la cláusula de escape en junio de 2025 marcó un punto de inflexión en la gestión fiscal del país. Amparado en la imposibilidad de cumplir con las obligaciones legales de gasto a partir de una proyección realista de ingresos, el Gobierno acudió al mecanismo previsto en la \href{https://www.funcionpublica.gov.co/eva/gestornormativo/norma.php?i=42956}{Ley 1473 de 2011}, que permite suspender temporalmente los límites al déficit establecidos por la regla fiscal, siempre y cuando se adopte una senda de retorno hacia la sostenibilidad de la deuda pública.

Dada la magnitud de las inflexibilidades legales del gasto y la ausencia de un marco institucional que permita resolver el conflicto entre dichas obligaciones y las restricciones de la regla fiscal, la activación de esta cláusula era, en la práctica, la única salida viable. La estrategia previa, basada en la sobreestimación sistemática de ingresos, ya se había agotado y había comprometido seriamente la credibilidad del marco fiscal.

Sin embargo, el uso de esta herramienta no está exento de riesgos. Existen temores legítimos de que, en ausencia de controles efectivos, la activación de la cláusula de escape se convierta en una excusa para relajar la disciplina fiscal y permitir que el gasto se desborde. Para evitarlo, es indispensable que el Gobierno actúe con transparencia, adopte lo más pronto posible un plan de consolidación creíble y comience a implementar medidas concretas para corregir los desequilibrios fiscales.

Esto implica apoyar las iniciativas legislativas orientadas a aumentar los ingresos y contener o reducir el gasto. La cláusula de escape no debe verse como una licencia para posponer decisiones difíciles, sino como una válvula institucional que ofrece un margen temporal para iniciar, sin dilaciones, un ajuste fiscal ordenado. En todo caso, será indispensable que el país rediseñe su marco institucional para que, agotadas otras alternativas, el Gobierno tenga la facultad de decretar los impuestos necesarios para cumplir con las obligaciones legales de gasto y respetar la regla fiscal.

\section{Acceso público al detalle del cierre fiscal}

La evaluación objetiva del cumplimiento de las metas fiscales requiere que organismos independientes —como el \gls{carf} o la \gls{oatp}— están en capacidad de validar las cifras fiscales reportadas por el Ministerio de Hacienda. Esto implica no solo el acceso al detalle de la ejecución presupuestal, sino también al mecanismo mediante el cual las cifras de ejecución presupuestal se convierten en las cifras del cierre fiscal. La necesidad de este control se hace evidente al examinar episodios recientes de contabilidad creativa. En 2021, por ejemplo, el déficit fiscal reportado fue del 7,0\% del PIB. Sin embargo, si se excluyera correctamente la enajenación de ISA —una operación de venta de activos que aumentó artificialmente los llamados ``recursos de capital'' en alrededor de un punto del PIB— el déficit habría sido de alrededor del 8,0\%. Y si, adicionalmente, se hubiera registrado el déficit del \gls{fepc} en el año en que efectivamente se causó, en lugar de en el año en que se pagó, el déficit real habría ascendido a una cifra cercana al 9,5\% del PIB. Estos ajustes ilustran cómo, al ser juez y parte, el Ministerio de Hacienda puede presentar una versión maquillada de su propio desempeño fiscal. Por ello, resulta urgente que exite acceso público al detalle sobre cómo la ejecución presupuestal se traduce en el cierre fiscal, de manera que cualquier interesado pueda auditar estas cifras, validar los supuestos y asegurar que las estadísticas fiscales reflejen de manera veraz la situación de las finanzas públicas.

\section{Conclusión}

Las propuestas aquí presentadas deben entenderse como una base para la discusión y el diseño de una agenda de reformas más amplia. No buscan reemplazar decisiones fiscales difíciles, sino complementar ese esfuerzo, dotando a diferentes ramas del poder público de herramientas más eficaces para gestionar el presupuesto en un contexto de crecientes inflexibilidades legales y constitucionales. Su objetivo es permitir que el ajuste entre ingresos y gastos se logre sin incumplir las obligaciones que el propio ordenamiento jurídico impone y sin comprometer la credibilidad del marco fiscal de mediano plazo.

Las reformas institucionales y administrativas son necesarias para mejorar la gestión de las finanzas públicas, fortalecer la disciplina fiscal y otorgarle al Ejecutivo mayores facultades para ajustar ingresos y gastos a lo largo del ciclo presupuestal, con la validación de instituciones independientes y bajo criterios más exigentes de transparencia. Sin embargo, por sí solas, estas reformas son insuficientes para resolver el desequilibrio estructural entre ingresos y gastos que enfrenta el país. Sin decisiones fiscales de fondo —que impliquen tanto medidas de contención y reducción del gasto como aumentos sostenidos de ingresos tributarios—, los avances institucionales corren el riesgo de convertirse en esfuerzos formales sin impacto real sobre la sostenibilidad de las cuentas públicas.